\documentclass[
    a4paper,
    chapterprefix,
    DIV=9
]{scrartcl}

%=========================================
%=============== Packages ================
% Fonts
\usepackage[utf8]{inputenc}
\usepackage[T1]{fontenc}
\usepackage{lmodern}
\usepackage{tgpagella}
\usepackage{mathpazo}

% Math formatting
\usepackage{mathtools}
\usepackage{amsfonts}
\usepackage{amsmath}
\usepackage{amssymb}
\usepackage{amsthm}
\usepackage{thmtools}

% Page formatting, colour and images
\usepackage{geometry}
\usepackage{graphicx}
\usepackage{tocloft}
\usepackage[x11names]{xcolor}
\usepackage{wrapfig}
\usepackage{cutwin}
\usepackage{tikz}
\usepackage{fancyhdr}
\usepackage{emptypage}
\usepackage{imakeidx}
\usepackage{mdframed}
\usepackage{eso-pic}
\usepackage{float}

% Hyperlinks and references
\usepackage{hyperref}
\usepackage[capitalise, nameinlink]{cleveref}
\usepackage{crossreftools}
\usepackage[
    backend=bibtex,
    style=alphabetic,
    sorting=ynt
]{biblatex}
\usepackage[
    totoc,
    columnsep=20pt,
    hangindent=8pt,
    subindent=20pt,
    subsubindent=30pt
]{idxlayout}
\usepackage[record,symbols]{glossaries-extra}

% Miscellaneous
\usepackage{multicol}
\usepackage{multirow}
\usepackage{enumitem}

%=========== (Re)definitions =============
% Symbols redefinitions
\let\oldemptyset\emptyset
\let\emptyset\varnothing

\let\totient\varphi

\renewcommand{\vert}{ \;\: \vline \;\: }
\newcommand{\vertalt}{ \;\: | \;\: }

%============= Formatting ================
% Line format
\linespread{1.05}

% Indent
\setlength{\parindent}{0pt}
\newlength{\normalparindent}
\setlength{\normalparindent}{\parindent}

% Heading format
\addtokomafont{disposition}{\rmfamily}

%============ Environments ===============
% Questions environments
\newenvironment{partquestions}[1]
{\begin{enumerate}[label=\textbf{(#1)}]}
{\end{enumerate}}

%=========== Custom Commands =============
\newcommand{\ZeroRoman}[1]{\ifcase\value{#1}\relax0\else\Roman{#1}\fi}  % Roman numeral that works with zero
\newcommand{\code}[1]{\texttt{#1}}  % Code block

%=========== General Symbols =============
\newcommand{\lcm}{\mathrm{lcm}}  % Lowest common multiple function
\newcommand{\sgn}{\mathrm{sgn}}  % Signum function

\newcommand{\im}{\mathrm{im}\;}  % Image of a function
\newcommand{\id}{\mathrm{id}}    % Identity function

\newcommand{\Z}{\mathbb{Z}}  % The set of integers
\newcommand{\Q}{\mathbb{Q}}  % The set of rational numbers
\newcommand{\R}{\mathbb{R}}  % The set of real numbers
\newcommand{\C}{\mathbb{C}}  % The set of complex numbers

%================ Part I =================
\newcommand{\An}[1]{\mathrm{A}_{#1}}                  % Alternating group of degree n
\newcommand{\Aut}[1]{\mathrm{Aut}(#1)}                % Group of automorphisms of G
\newcommand{\Centralizer}[2]{\mathrm{C}_{#1}(#2)}     % Centralizer of an element in G
\newcommand{\Cl}[1]{\mathrm{Cl}(#1)}                  % Conjugacy class of the element x
\newcommand{\Cn}[1]{\mathrm{C}_{#1}}                  % Cyclic group of order n
\newcommand{\GL}[2]{\mathrm{GL}_{#1}\left(#2\right)}  % General Linear Group of degree n
\newcommand{\Inn}[1]{\mathrm{Inn}(#1)}                % Group of inner automorphisms of G
\newcommand{\N}[2]{\mathrm{N}_{#1}(#2)}               % Normalizer of S in G
\newcommand{\Out}[1]{\mathrm{Out}(#1)}                % Group of outer automorphisms of G
\newcommand{\SL}[2]{\mathrm{SL}_{#1}\left(#2\right)}  % Special Linear Group of degree n
\newcommand{\Sn}[1]{\mathrm{S}_{#1}}                  % Symmetric group of degree n
\newcommand{\Syl}[2]{\mathrm{Syl}_{#1}(#2)}           % Set of Sylow p-groups of G
\newcommand{\Sym}[1]{\mathrm{Sym}(#1)}                % Symmetric group of a set
\newcommand{\Un}[1]{\mathcal{U}_{#1}}                 % Group of units modulo n
\newcommand{\CenterGrp}[1]{\mathrm{Z}(#1)}            % Center of a group G

\newcommand{\Stab}[2]{\mathrm{Stab}_{#1}(#2)}         % Stabilizer of x by G
\newcommand{\Fix}[2]{\mathrm{Fix}_{#1}(#2)}           % Set of all elements in X which is fixed by g
\newcommand{\Orb}[2]{\mathrm{Orb}_{#1}(#2)}           % Orbit of x under G

%================ Part II ================
\newcommand{\ideal}[1]{\mathfrak{#1}}                % Ideal of a ring
\newcommand{\princ}[1]{\left\langle#1\right\rangle}  % Principal ideal generated by the element

\newcommand{\Mn}[2]{\mathcal{M}_{#1\times#1}(#2)}  % The ring of n by n matrices with entries in the ring R
\newcommand{\ZeroM}[1]{\textbf{0}_{#1}}      % Zero matrix
\newcommand{\IdentityM}[1]{\textbf{I}_{#1}}  % Identity matrix

% \renewcommand{\H}{\mathbb{H}}   % Quaternion Ring
% \newcommand{\qi}{\textbf{i}}    % Quaternion i
% \newcommand{\qj}{\textbf{j}}    % Quaternion j
% \newcommand{\qk}{\textbf{k}}    % Quaternion k

\newcommand{\Ann}[2]{\mathrm{Ann}_{#1}(#2)}  % Annihilator of a subset A of a ring R
\newcommand{\Char}[1]{\mathrm{char}(#1)}     % Characteristic of a ring R
\newcommand{\Nilr}[1]{\mathfrak{N}_{#1}}     % Nilradical of a ring R
\newcommand{\Frac}[1]{\mathrm{Frac}(#1)}     % Field of fractions of an integral domain D

%=============== Part III ================
\newcommand{\Span}[1]{\mathrm{span}(#1)}  % Span of a set S of vectors
\let\olddim\dim
\renewcommand*{\dim}[1]{\olddim(#1)}      % Dimension of a vector space V


%============ Other Setup ================
\hypersetup{
    pdfborder = {0 0 0}
}


%=========================================
\title{\bfseries\scshape A Complete Introduction To\\Abstract Algebra}
\subtitle{\bfseries\scshape Errata}
\date{\texttt{2026-06-01 1940H}}
\author{Kan Onn Kit}

\begin{document}
\maketitle

Despite my best efforts, some errors have unfortunately been made in \textit{A Complete Introduction To Abstract Algebra}. This document lists the errors that have been made, and their corrections.

\section*{In the 1\textsuperscript{st} Edition (Post 2) [and Before]}
\subsection*{Major Errors}
\begin{itemize}
    \item The subgroup of $\mathrm{C}_2$ in Example 18.3.3 is wrong: the generator should be $r^3$ instead of $s$. So instead of ``$\{e, \underline{s}\} = \langle \underline{s} \rangle$'' it should have been ``$\{e, \underline{r^3}\} = \langle \underline{r^3} \rangle$''
\end{itemize}

\subsection*{Minor Errors}
\begin{itemize}
    \item There is a minor error in the statement after Figure 13.2's caption (Subgroup Lattice for Theorem 13.4.1), where the original text mistakenly omits the condition for $H$ being a normal subgroup of $G$ for almost all the conditions except for the ``dotted quotient groups''. In particular, the section
    \begin{quote}
        We only show subgroups that we care about in the figure above. The group $G$ has \textit{..[truncated]..} subgroup $H\cap K$. The dotted quotient groups are isomorphic to each other if $H \unlhd G$.
    \end{quote}
    should have read
    \begin{quote}
        We only show subgroups that we care about in the figure above. So, if $H$ is a normal subgroup of $G$, then $G$ has \textit{..[truncated]..} subgroup $H\cap K$; and the dotted quotient groups are isomorphic to each other.
    \end{quote}
    \item In Example 18.3.3, the subgroup of $D_3$ could also be $\{e, r^2, r^4, rs, r^3s, r^5s\}$. The full text should thus read ``$D_3$ is $\{e, r^2, r^4, s, r^2s, r^4s\}$ \underline{(or $\{e, r^2, r^4, rs, r^3s, r^5s\}$)}''.
\end{itemize}

\subsection*{Typos}
None found.

\section*{In the 1\textsuperscript{st} Edition (Post 1) [and Before]}
\subsection*{Major Errors}
None found.

\subsection*{Minor Errors}
None found.

\subsection*{Typos}
\begin{itemize}
    \item There is a minor grammatical error in the proof of lemma 9.4.3, where the phrase ``second bullet point \underline{is} proven by'' should have been ``second bullet point \underline{was} proven by''.
    \item There is a minor grammatical error in the proof of lemma 9.4.3, where the phrase ``$H$ and $K$ \underline{has} a common subgroup'' to ``$H$ and $K$ \underline{have} a common subgroup''.
    \item In the list of frequently used notation,
    \begin{itemize}
        \item ``\underline{Stabilzer} of $x$ by $G$'' should have been ``\underline{Stabil\textbf{\textit{i}}zer} of $x$ by $G$''; and
        \item ``Principal ideal \underline{generate\textbf{\textit{e}}d} by'' should have been ``Principal ideal \underline{generated} by''.
    \end{itemize}
\end{itemize}

\section*{In the 1\textsuperscript{st} Edition [and Before]}
\subsection*{Major Errors}
None found.

\subsection*{Minor Errors}
\begin{itemize}
    \item In example 5.1.5, I missed out a few elements in the set $R$. The full set should be
    \begin{align*}
        R &= \{(3, 3), (3, 5), (3, 7), (4, 4), (4, 6), (5, 3), (5, 5),\\
        &\quad\quad(5, 7), (6, 4), (6, 6), (7, 3), (7, 5), (7, 7)\}.
    \end{align*}
\end{itemize}

\subsection*{Typos}
\begin{itemize}
    \item There is a small typo in the answer of exercise 1.1(f). The non-italicised "S" should be italicised (as it is the name of a set).
\end{itemize}

\begin{center}
    \vfill
    {\Large\bfseries\scshape End of Errata}
\end{center}
\end{document}
